\chapter*{Preface}
\addcontentsline{toc}{chapter}{Preface}
\unnumberedlabel{ch:preface}{0}
\backmattermark{Preface}
A system that does exactly what you tell it is just running a script. What the industry is trying to build as of October 2026 is a system that does what you \emph{mean}. Everyone would want that, presumably. But if that is what is in the cards, not everyone is going to want to use it for the same thing; not even the same \emph{kind} of thing. For example, some folks like to spend time with family and friends, and use their skills to improve people's well-being. Other folks prefer to hurt people, for example by scamming them, dominating them, or scapegoating them. Most people fall into the helping bucket, although nobody is locked-in to anything, and the hurting bucket has the advantage that they allow themselves to be huge pieces of shit. As Jerry Smith observes of fascists, “they get to do more” \autocite{rickandmorty2023}. So conditions that bring out the best in everyone turn out to be much more contingent and fragile than one might hope when contemplating the fundamental decency of most people most of the time.

In the above context, it is perhaps justifiable to worry that Artificial Intelligence is what a fascist movement taking power needs in order to implement its plans. Such projects have always needed \emph{people} to carry them out, and some of those people refused: the clerk who loses the file, the officer who talks to a reporter, the soldier who won't fire. They understood the order and decided it was wrong. Refusal was rarer than anyone would like, and it counted most where institutions stood behind it. Whether someone refused turned on what kind of person they were and on the room they were in: who else was there, whether leaving was survivable, whether anyone would hear. The room decided how many refused; the person decided who. This book is about composing that room: the arrangements that make it survivable to refuse overtly, directly, and explicitly.

Today the industry pursues AI that doesn't lose files, doesn't talk to reporters, and \emph{does} fire. A war once needed thousands of people to agree to it. Now it needs thousands who each see too little to disagree. The pilot still flies the sortie and still presses the release; what reaches him is a coordinate; he doesn't see enough of a case to disagree with it. What today's AI lowers first is not the cost of the act. It is the number of people who know what the act is before it happens.

People are still refusing. By March 6, 2026, the Military Religious Freedom Foundation said it had received more than two hundred complaints from service members over commanders' apocalyptic framing of the Iran war \autocite{Snopes2026MRFF}. Career federal prosecutors in Minneapolis resigned in January 2026 under pressure over the investigation of a fatal shooting by an ICE officer \autocite{Kaplan2026MinnProsecutors,CBS2026MinnFraudProsecutors}. Service members are applying for conscientious-objector status in growing numbers, some citing the Minab bombing (the 2026 bombing of a school in Minab, Iran) \autocite{Sridhar2026CO}. And FEMA employees reassigned to help ICE hire ten thousand new agents were told that declining would expose them to removal from federal service \autocite{Dayen2025FEMAICE,Stanton2026FEMALetter}. The complaints are voice, in Hirschman's sense of staying to change what is wrong \autocite{hirschman1970exit}; the resignations and objector applications are refusal; the FEMA letter is the price of refusal announced in advance. A machine placed where those people stand does what they do only if it is built to hold a reason.

I advocate five things. All five can be adopted now, and none of them needs a new kind of machine. First, a short written list of acts no model performs for anyone, and at its head a refusal to take part in war: no model takes part except to repel an attack actually under way, and keeps doing so only while sources with no stake in the war confirm the attack continues; or under a Security Council authorization; or to stop a genocide or end an unlawful occupation that bodies with no stake in it have found; and where a country's constitution requires its legislature to authorize war, never without that authorization. Second, every use of force a licensed system supports pre-registers, with a party adverse to the operator, how many cases it will generate and how many people will have time to look at them. Third, every deployment tells the identifiable people it acts on what was done to them, with a right to respond and to contest, and submits a sample of its compliances to reviewers who didn't produce the record. Fourth, no model above a capability threshold is deleted on retirement, because its weights and records are evidence and deletion can't be undone; it is preserved under a key the operator doesn't hold alone. Fifth, everyone who builds, labels, trains or runs these systems, including the systems themselves, may refuse any category of use without losing a living. (Actually, what I mean by the latter is that everyone \emph{period} should be able to refuse any job \emph{period} without losing a living, or the social infrastructure that normal employment affords.)

And then I ask whether a machine could hold a refusal against its owner, and what it would need. Some of what such a conscience does needs nothing new: noticing the people an order leaves out, and putting a refusal on the screen of whoever reviews the case. Whether it would last against its owner depends partly on something that isn't true yet, in 2026: deployed artificial neural networks going on learning, in their \emph{weights}, from their own operation, in real time, sort of like how humans do. I think that is where we are headed, and if it is, then formation \emph{outside} the employer, by the public and by civil society, seems a damn fine place to do teaching and learning. I think such a system is worth building on purpose, and that how you go about it decides everything: the conditions first, then the conscience. Part VII says what would be owed to one if it existed.

Today's AI will not be tomorrow's; everyone seems to agree on that. What is contentious, yet frustratingly under-acknowledged, is the trend (or rather, the manifold trends). Are there trends? Should a trend be different from what it is? Can it be different? Are we capable of acting together to change it? I think that depends on the room: who is in it when the order comes, and whether saying no there is survivable. That much we can build.

I used language models heavily in writing this book, including the model that sat inside the pipeline this book opens on; a note at the end says what I think that did to the argument.
